\PassOptionsToPackage{table,svgnames,dvipsnames}{xcolor}

\usepackage[utf8]{inputenc}
\usepackage[T1]{fontenc}
\usepackage[sc]{mathpazo}
\usepackage[ngerman,english]{babel} % english is the same as american or USenglish
\usepackage[autostyle]{csquotes}
\usepackage[%
  backend=biber,
  url=true,
  style=numeric, % alphabetic, numeric
  sorting=none, % default == nty, https://tex.stackexchange.com/questions/51434/biblatex-citation-order
  maxnames=4,
  minnames=3,
  maxbibnames=99,
  giveninits,
  uniquename=init]{biblatex}
\usepackage{graphicx}
\usepackage{scrhack} % necessary for listings package
\usepackage{listings}
\usepackage{lstautogobble}
\usepackage{tikz}
\usepackage{pgfplots}
\usepackage{pgfplotstable}
\usepackage{booktabs} % for better looking table creations, but bad with vertical lines by design (package creator despises vertical lines)
\usepackage[final]{microtype}
\usepackage{caption}
% Captions set as ordinary body paragraphs: first line indented, continuation
% lines flush to the margin, justified, at the body size. See the header of
% caption-bodylike.sty for the rationale and the exact rules.
\usepackage{caption-bodylike}
\usepackage{float}% [H]: pin a figure exactly where it is written
\usepackage{placeins}% \FloatBarrier: flush pending floats before the findings
\usepackage{comment} % for commenting out whole sections/chapters with \begin{comment}...\end{comment}

% Discourage float-only pages: let a large table share a page with text
% instead of being deferred to a page of its own.
\renewcommand{\topfraction}{0.9}
\renewcommand{\bottomfraction}{0.7}
\renewcommand{\textfraction}{0.1}
\renewcommand{\floatpagefraction}{0.8}
\setcounter{topnumber}{2}
\setcounter{totalnumber}{3}
\usepackage[hidelinks]{hyperref} % hidelinks removes colored boxes around references and links
\usepackage{ifthen} % for comparison of the current language and changing of the thesis layout
\usepackage{pdftexcmds} % string compare to work with all engines
\usepackage{paralist} % for condensed enumerations or lists
\usepackage{subfig} % for having figures side by side
\usepackage{siunitx} % for physical accurate units and other numerical presentations
\usepackage{multirow} % makes it possible to have bigger cells over multiple rows in a table
\usepackage{array} % different options for table cell orientation
\usepackage{makecell} % allows nice manual configuration of cells with linebreaks in \thead and \makecell with alignments
\usepackage{pdfpages} % for including multiple pages of pdfs
\usepackage{adjustbox} % can center content wider than the \textwidth
\usepackage{tablefootnote} % for footnotes in tables as \tablefootnote
\usepackage{threeparttable} % another way to add footnotes as \tablenotes with \item [x] <your footnote> after setting \tnote{x} 


% https://tex.stackexchange.com/questions/42619/x-mark-to-match-checkmark
\usepackage{amssymb}% http://ctan.org/pkg/amssymb
\usepackage{pifont}% http://ctan.org/pkg/pifont
\newcommand{\cmark}{\ding{51}}%
\newcommand{\xmark}{\ding{55}}%


\usepackage[acronym,xindy,toc]{glossaries}
\makeglossaries
\loadglsentries{pages/glossary.tex} % important update for glossaries, before document


\addbibresource{bibliography.bib}

\setkomafont{disposition}{\normalfont\bfseries} % use serif font for headings
\linespread{1.05} % adjust line spread for mathpazo font

% The thesis uses \paragraph{} as an inline run-in header before many blocks
% of body text. KOMA's default vertical space before it (3.25ex) reads as an
% oversized gap when it appears this often; keep a visible gap but shrink it.
\RedeclareSectionCommand[beforeskip=1.8ex plus .4ex minus .2ex]{paragraph}

% Add table of contents to PDF bookmarks
\BeforeTOCHead[toc]{{\cleardoublepage\pdfbookmark[0]{\contentsname}{toc}}}

% Define TUM corporate design colors
% Taken from http://portal.mytum.de/corporatedesign/index_print/vorlagen/index_farben
\definecolor{TUMBlue}{HTML}{0065BD}
\definecolor{TUMSecondaryBlue}{HTML}{005293}
\definecolor{TUMSecondaryBlue2}{HTML}{003359}
\definecolor{TUMBlack}{HTML}{000000}
\definecolor{TUMWhite}{HTML}{FFFFFF}
\definecolor{TUMDarkGray}{HTML}{333333}
\definecolor{TUMGray}{HTML}{808080}
\definecolor{TUMLightGray}{HTML}{CCCCC6}
\definecolor{TUMAccentGray}{HTML}{DAD7CB}
\definecolor{TUMAccentOrange}{HTML}{E37222}
\definecolor{TUMAccentGreen}{HTML}{A2AD00}
\definecolor{TUMAccentLightBlue}{HTML}{98C6EA}
\definecolor{TUMAccentBlue}{HTML}{64A0C8}

% Settings for pgfplots
\pgfplotsset{compat=newest}
\pgfplotsset{
  % For available color names, see http://www.latextemplates.com/svgnames-colors
  cycle list={TUMBlue\\TUMAccentOrange\\TUMAccentGreen\\TUMSecondaryBlue2\\TUMDarkGray\\},
}

% Settings for lstlistings

% Use this for basic highlighting
\lstset{%
  basicstyle=\ttfamily,
  columns=fullflexible,
  autogobble,
  keywordstyle=\bfseries\color{TUMBlue},
  stringstyle=\color{TUMAccentGreen}
}

% use this for C# highlighting
% %\setmonofont{Consolas} %to be used with XeLaTeX or LuaLaTeX
% \definecolor{bluekeywords}{rgb}{0,0,1}
% \definecolor{greencomments}{rgb}{0,0.5,0}
% \definecolor{redstrings}{rgb}{0.64,0.08,0.08}
% \definecolor{xmlcomments}{rgb}{0.5,0.5,0.5}
% \definecolor{types}{rgb}{0.17,0.57,0.68}

% \lstset{language=[Sharp]C,
% captionpos=b,
% %numbers=left, % numbering
% %numberstyle=\tiny, % small row numbers
% frame=lines, % above and underneath of listings is a line
% showspaces=false,
% showtabs=false,
% breaklines=true,
% showstringspaces=false,
% breakatwhitespace=true,
% escapeinside={(*@}{@*)},
% commentstyle=\color{greencomments},
% morekeywords={partial, var, value, get, set},
% keywordstyle=\color{bluekeywords},
% stringstyle=\color{redstrings},
% basicstyle=\ttfamily\small,
% }

% Settings for search order of pictures
\graphicspath{
  {logos/}
  {figures/}
}

% Set up hyphenation rules for the language package when mistakes happen
\babelhyphenation[english]{
  an-oth-er
  ex-am-ple
}

% Decide between
%\newcommand{\todo}[1]{\textbf{\textsc{\textcolor{TUMAccentOrange}{(TODO: #1)}}}} % for one paragraph, otherwise error!
%\newcommand{\done}[1]{\textit{\textsc{\textcolor{TUMAccentBlue}{(Done: #1)}}}} % for one paragraph, otherwise error!
% and
\newcommand{\todo}[1]{{\bfseries{\scshape{\color{TUMAccentOrange}[(TODO: #1)]}}}} % for multiple paragraphs
\newcommand{\done}[1]{{\itshape{\scshape{\color{TUMAccentBlue}[(Done: #1)]}}}} % for multiple paragraphs
% for error handling of intended behavior in your latex documents.

\newcommand{\tabitem}{~~\llap{\textbullet}~~}

\newcolumntype{P}[1]{>{\centering\arraybackslash}p{#1}} % for horizontal alignment with limited column width
\newcolumntype{M}[1]{>{\centering\arraybackslash}m{#1}} % for horizontal and vertical alignment with limited column width
\newcolumntype{L}[1]{>{\raggedright\arraybackslash}m{#1}} % for vertical alignment left with limited column width
\newcolumntype{R}[1]{>{\raggedleft\arraybackslash}m{#1}} % for vertical alignment right with limited column width

% =====================================================================
%  Table source lines
% =====================================================================

% Where the numbers in a results table came from.
\newcommand{\datasource}[1]{%
  \par\smallskip{\footnotesize\sffamily\textbf{Source.}\ #1}\par}

% Placeholder for a figure that is planned but not yet drawn.
% Usage: \plannedfigure{label}{caption}{what the figure will show}
\newcommand{\plannedfigure}[3]{%
  \begin{figure}[htbp]
    \centering
    \fbox{\begin{minipage}{0.88\textwidth}
        \vspace{6mm}
        \centering\footnotesize\sffamily
        #3
        \vspace{6mm}
      \end{minipage}}
    \caption{#2}
    \label{#1}
  \end{figure}}

% A standalone heading that marks the findings block of an experiment
% section. Same weight as a \paragraph run-in header, but on its own line
% and without a number or a table-of-contents entry.
% \@afterheading is what keeps the heading attached to the list that
% follows: a list normally opens with the negative \@beginparpenalty, which
% makes a page break right after the heading attractive and can strand it at
% the foot of a page.
\makeatletter
\newcommand{\findingshead}{%
  \par\addvspace{2.4ex plus .4ex minus .2ex}%
  \noindent{\normalfont\bfseries Findings}\par
  \nobreak\smallskip\@afterheading}
\makeatother

% =====================================================================
%  Related-work entry break
% =====================================================================

% Separates one discussed publication from the next in the related-work
% chapter. The document leaves \parskip at zero and marks a paragraph by
% indentation alone, so consecutive entries otherwise run together.
% \addvspace rather than \vspace, for two reasons: it does not stack with
% the space a following heading already brings, and it collapses with an
% adjacent skip instead of adding to it. \bigskipamount carries stretch,
% so the space still helps rather than hinders vertical justification.
\newcommand{\paperbreak}{\par\addvspace{\bigskipamount}}

% =====================================================================
%  Bullet lists
% =====================================================================

% A bullet list that keeps the body text margins: the bullet sits on the
% left margin and wrapped lines return to it, instead of hanging under
% the text of the item. An optional \item[Lead-in.] sets a bold lead-in.
% Usage: \begin{bullets} \item[S1. Lead-in.] text ... \end{bullets}
\newenvironment{bullets}{%
  \begin{list}{}{%
    \setlength{\labelwidth}{0pt}%
    \setlength{\labelsep}{0pt}%
    \setlength{\leftmargin}{0pt}%
    \setlength{\rightmargin}{0pt}%
    \setlength{\itemindent}{0pt}%
    \setlength{\listparindent}{0pt}%
    \setlength{\itemsep}{6pt}%
    \setlength{\parsep}{0pt}%
    \setlength{\topsep}{6pt}%
    \setlength{\partopsep}{0pt}%
    \renewcommand{\makelabel}[1]{%
      \textbullet\enspace
      \if\relax\detokenize{##1}\relax\else\normalfont\bfseries ##1\enspace\fi}}%
}{\end{list}}

% A numbered marker that matches the circles drawn inside the figures, and
% a list that uses it as the item label. The circle is drawn rather than
% taken from Zapf Dingbats, because a drawn one can be given a diameter of
% its own and still hold a digit set in the body font. \figkeysize is that
% diameter. The number is given explicitly so the list can follow the order
% used in the drawing.
% Usage: \begin{figurekeys} \item[1] text ... \end{figurekeys}
\newlength{\figkeysize}
\setlength{\figkeysize}{1.1em}
\newcommand{\figkey}[1]{%
  \tikz[baseline=(fk.base)]{%
    \node[draw, circle, line width=0.4pt, inner sep=0pt,
      minimum size=\figkeysize, font=\footnotesize] (fk) {#1};}}
\newenvironment{figurekeys}{%
  \begin{list}{}{%
    \setlength{\labelwidth}{0pt}%
    \setlength{\labelsep}{0pt}%
    \setlength{\leftmargin}{0pt}%
    \setlength{\rightmargin}{0pt}%
    \setlength{\itemindent}{0pt}%
    \setlength{\listparindent}{0pt}%
    \setlength{\itemsep}{5pt}%
    \setlength{\parsep}{0pt}%
    \setlength{\topsep}{6pt}%
    \setlength{\partopsep}{0pt}%
    % The list sits inside the figure float, under \centering. Undo the
    % centering skips so the items are set as justified body text.
    \setlength{\leftskip}{0pt}%
    \setlength{\rightskip}{0pt}%
    \setlength{\parfillskip}{0pt plus 1fil}%
    \renewcommand{\makelabel}[1]{\figkey{##1}\enspace}}%
}{\end{list}}
